%%%%%%%%%%%%%%%%%%%%%%%%%%%%%%%%%%%%%%%%%
% Short Sectioned Assignment LaTeX Template Version 1.0 (5/5/12)
% This template has been downloaded from: http://www.LaTeXTemplates.com
% Original author:  Frits Wenneker (http://www.howtotex.com)
% License: CC BY-NC-SA 3.0 (http://creativecommons.org/licenses/by-nc-sa/3.0/)
%%%%%%%%%%%%%%%%%%%%%%%%%%%%%%%%%%%%%%%%%

% \documentclass[paper=a4, fontsize=11pt]{scrartcl} % A4 paper and 11pt font size
\documentclass[11pt, a4paper]{book}
\usepackage[T1]{fontenc} % Use 8-bit encoding that has 256 glyphs
\usepackage[utf8]{inputenc}
\usepackage{fourier} % Use the Adobe Utopia font for the document - comment this line to return to the LaTeX default
\usepackage{listings} % para insertar código con formato similar al editor
\usepackage[spanish, es-tabla]{babel} % Selecciona el español para palabras introducidas automáticamente, p.ej. "septiembre" en la fecha y especifica que se use la palabra Tabla en vez de Cuadro
\usepackage{url} % ,href} %para incluir URLs e hipervínculos dentro del texto (aunque hay que instalar href)
\usepackage{graphics,graphicx, float} %para incluir imágenes y colocarlas
\usepackage[gen]{eurosym} %para incluir el símbolo del euro
\usepackage{cite} %para incluir citas del archivo <nombre>.bib
\usepackage{enumerate}
\usepackage{hyperref}
\usepackage{graphicx}
\usepackage{tabularx}
\usepackage{booktabs}

\usepackage[table,xcdraw]{xcolor}
\hypersetup{
	colorlinks=true,	% false: boxed links; true: colored links
	linkcolor=black,	% color of internal links
	urlcolor=cyan		% color of external links
}
\renewcommand{\familydefault}{\sfdefault}
\usepackage{fancyhdr} % Custom headers and footers
\pagestyle{fancyplain} % Makes all pages in the document conform to the custom headers and footers
\fancyhead[L]{} % Empty left header
\fancyhead[C]{} % Empty center header
\fancyhead[R]{Me Llamo Así} % My name
\fancyfoot[L]{} % Empty left footer
\fancyfoot[C]{} % Empty center footer
\fancyfoot[R]{\thepage} % Page numbering for right footer
%\renewcommand{\headrulewidth}{0pt} % Remove header underlines
\renewcommand{\footrulewidth}{0pt} % Remove footer underlines
\setlength{\headheight}{13.6pt} % Customize the height of the header

\usepackage{titlesec, blindtext, color}
\definecolor{gray75}{gray}{0.75}
\newcommand{\hsp}{\hspace{20pt}}
\titleformat{\chapter}[hang]{\Huge\bfseries}{\thechapter\hsp\textcolor{gray75}{|}\hsp}{0pt}{\Huge\bfseries}
\setcounter{secnumdepth}{4}
\usepackage[Lenny]{fncychap}


\begin{document}

	% Plantilla portada UGR
	\begin{titlepage}
\newlength{\centeroffset}
\setlength{\centeroffset}{-0.5\oddsidemargin}
\addtolength{\centeroffset}{0.5\evensidemargin}
\thispagestyle{empty}

\noindent\hspace*{\centeroffset}\begin{minipage}{\textwidth}

\centering
\includegraphics[width=0.9\textwidth]{logos/logo_ugr.jpg}\\[1.4cm]

\textsc{ \Large TRABAJO FIN DE GRADO\\[0.2cm]}
\textsc{ GRADO EN INGENIERIA INFORMATICA}\\[1cm]

{\Huge\bfseries Título \\}
\noindent\rule[-1ex]{\textwidth}{3pt}\\[3.5ex]
{\large\bfseries Subtítulo }
\end{minipage}

\vspace{2.5cm}
\noindent\hspace*{\centeroffset}
\begin{minipage}{\textwidth}
\centering

\textbf{Autor}\\ {Estudiante}\\[2.5ex]
\textbf{Director}\\ {Tutor(a)(es)}\\[2cm]
\includegraphics[width=0.3\textwidth]{logos/etsiit_logo.png}\\[0.1cm]
\textsc{Escuela Técnica Superior de Ingenierías Informática y de Telecomunicación}\\
\textsc{---}\\
Granada, Junio de 201x
\end{minipage}
\end{titlepage}


	% Plantilla prefacio UGR
	\thispagestyle{empty}

\begin{center}
{\large\bfseries Héroes \\ Geolocalización para ayuda en emergencias }\\
\end{center}
\begin{center}
Alejandro Menor Molinero\\
\end{center}

%\vspace{0.7cm}

\vspace{0.5cm}

\noindent{\textbf{Resumen}\\

La inseguridad en la vía pública, especialmente de las grandes ciudades, es un problema que muchas personas, sobre todo las mujeres, sufren cada día.
En mayor o menos medida nos podemos sentir vulnerables en la calle y nos gustaría estar acompañados. Héroes es una aplicación que permite conectar con personas
cercanas dispuestas a ayudar en muy poco tiempo. \\ \\
Para esto se ha implementado un servicio de geolocalización escalable con \mbox{NodeJS}, utilizando como infraestructura Redis y MongoDB, además de una aplicación
móvil con Flutter. El sistema permite alertar a personas cercanas, permitiéndolas seguir en tiempo real la ubicación de la víctima.
	

\cleardoublepage

\begin{center}
{\large\bfseries Héroes \\ Aid through geolocalization in emergencies}\\
\end{center}
\begin{center}
Alejandro Menor Molinero\\
\end{center}
\vspace{0.5cm}

\noindent{\textbf{Abstract}\\ \\
Insecurity in the streets, specially in big cities, it's a problem that many people, especially women, suffer every day.
We can feel vulnerable outside and we'd rather have company. Héroes is an app that allows us to connect with people nearby willing to help. \\ \\
To make this possible, an escalable geolocalization service has been implemented with NodeJS, using MongoDB and Redis as infraestructure, besides 
a mobile App with Flutter. The system sends an alert to nearby people, enabling them to follow the victim's location in real time.


\cleardoublepage

\thispagestyle{empty}

\noindent\rule[-1ex]{\textwidth}{2pt}\\[4.5ex]

D. \textbf{Juan Julián Merelo Guervós}, profesor del departamento de Arquitectura y Tecnología de Computadores.

\vspace{0.5cm}

\textbf{Informo:}

\vspace{0.5cm}

Que el presente trabajo, titulado \textit{\textbf{Héroes}},
ha sido realizado bajo mi supervisión por \textbf{Alejandro Menor Molinero}, y autorizo la defensa de dicho trabajo ante el tribunal
que corresponda.

\vspace{0.5cm}

Y para que conste, expiden y firman el presente informe en Granada a Julio de 2021.

\vspace{1cm}

\textbf{El/la director(a)/es: }

\vspace{5cm}

\noindent \textbf{Juan Julián Merelo Guervós}







	% Índice de contenidos
	\newpage
	\tableofcontents

	% Índice de imágenes y tablas
	\newpage
	\listoffigures

	% Si hay suficientes se incluirá dicho índice
	\listoftables 
	\newpage

	% Introducción 
	\chapter{Introducción}

Este proyecto es software libre, y está liberado con la licencia \cite{gplv3}.

\section{El problema}
En pleno 2021, en España, la calle no es siempre un lugar seguro. Dependiendo de la zona por la que te muevas, la hora o tu género
hay más o menos probabilidades de que te lleves, en el mejor de los casos, un susto.

Los últimos estudios, como el de RACC y Zurich \cite{racc-zurich}, sacan conclusiones muy preocupantes, por ejemplo: 
\begin{itemize}
  \item El 45\% de las mujeres encuestadas declara que ha sido víctima de acoso de noche en la vía pública yendo a pie.
  \item El 68\% y 37\% de mujeres y hombres respectivamente ha cambiado alguna vez de modo de transporte por motivos de seguridad personal.
\end{itemize}

\section{Estado del arte}\label{sec:art}
Durante los últimos años, han ido emergiendo distintas aplicaciones que tratan de paliar esta inseguridad que podemos clasificar en dos grupos:
\begin{enumerate}
\item \textbf{Botón del pánico a la policía:} En caso de estar en una situación complicada, podemos avisar a la policía sin necesidad de llamar. La más famosa es \textbf{AlertCops} que cuenta con el respaldo del ministerio del interior.
\item \textbf{Botón del pánico a un grupo de contactos:} Con una activación similar a el grupo anterior, el aviso esta vez se hace a una serie de personas cercanas. En muchos casos, estas serán capaces de saber donde estás, escuchar y ver lo que pasa a partir de ese momento. Algún ejemplo: \textbf{bSafe}, \textbf{When and Where} o \textbf{Companion}.
Otra ''implementación'', es la creación de grupos de mensajería masivos de estudiantes para pedir ayuda si se sienten inseguras.
\end{enumerate}

\section{Inconvenientes}
Pese a la cantidad de alternativas y a la alta calidad de muchas de ellas, no han conseguido tener el impacto que esperaban.
\begin{enumerate}
  \item La policía o los contactos están, en muchos casos, lejos de ti y la ayuda tiene que ser inmediata.
  \item A veces el peligro no es explícito y pedir auxilio es excesivo, por ejemplo: sospechas que alguien te está siguiendo.
\end{enumerate}

\section{Análisis del problema}
Queremos reducir al máximo la inseguridad de todas las personas en la vía pública.

Podemos identificar tres niveles dentro del problema:
\begin{enumerate}
  \item \textbf{Sospecha, ligera inseguridad:} No lo sabes con certeza, pero algo te pone alerta, podrías estar en peligro y te gustaría estar acompañado.
  \item \textbf{Miedo:} Estás en una situación totalmente incómoda, no sabes si va a llegar a más, pero necesitas compañía.
  \item \textbf{Pánico:} Es necesaria una ayuda inmediata y sería pertinente avisar a los cuerpos de seguridad.
\end{enumerate}

Para los dos primeros, que a su vez son las más comunes, no hay propuestas.

\section{Objetivos generales}\label{sec:obj}
\begin{enumerate}
  \item La solución tiene que cubrir la funcionalidad de los botones del pánico existentes, analizados en la sección \ref{sec:art}.
  Dar la posibilidad de: 
  \begin{itemize}
    \item Avisar a la policía con un botón en una situación de auxilio.
    \item Avisar a un grupo de contactos de confianza con un botón, compartiendo ubicación, audio y video.
  \end{itemize}
  Configurable por el usuario.
  \item Deberá recoger toda la información relevante del entorno que pueda servir para esclarecer los hechos.
  \item La solución tiene que ser efectiva para los primeros niveles de inseguridad descritos anteriormente.
  \begin{itemize}
    \item Dar la posibilidad de que alguien que esté cerca te acompañe si lo necesitas.
    \item Dar la posiblidad de prestar esa ayuda.
  \end{itemize}
En este trabajo, vamos a plantear y a implementar una solución para este objetivo.
\end{enumerate}



	% Descripción del problema y hasta donde se llega
	\chapter{Descripción del problema}



	% Estado del arte
	% 	1. Crítica al estado del arte
	% 	2. Propuesta
	\chapter{Estado del arte}
\section{Algunas propuestas}
Durante los últimos años, han sido emergiendo distintas aplicaciones que tratan de paliar esta inseguridad que podemos clasificar en dos grupos:
\begin{enumerate}
\item \textbf{Botón del pánico a la policía:} En caso de estar en una situación complicada, podemos avisar a la policía sin necesidad de llamar. La más famosa es \textbf{AlertCops} que cuenta con el respaldo del ministerio del interior.
\item \textbf{Botón del pánico a un grupo de contactos:} Con una activación similar a el grupo anterior, el aviso esta vez se hace a una serie de personas cercanas. En muchos casos, estas serán capaces de saber donde estás, escuchar y ver lo que pasa a partir de ese momento. Algún ejemplo: \textbf{bSafe}, \textbf{When and Where} o \textbf{Companion}.
Otra ''implementación'', es la creación de grupos de mensajería masivos de estudiantes para pedir ayuda si se sienten inseguras.
\end{enumerate}

\section{Inconvenientes}
Pese a la cantidad de alternativas y a la alta calidad de muchas de ellas, no han conseguido tener el impacto que esperaban.
\begin{enumerate}
  \item La policía o los contactos están, en muchos casos, lejos de ti y la ayuda tiene que ser inmediata.
  \item A veces el peligro no es explícito y pedir auxilio es excesivo, por ejemplo sospechas que alguien te está siguiendo.
\end{enumerate}


	
	\chapter{Planificación}

\section{Metodología utilizada}


\section{Temporización}

\section{Seguimiento del desarrollo}


	% Análisis del problema
	% 1. Análisis de requisitos
	% 2. Análisis de las soluciones
	% 3. Solucion propuesta
	% 4. Análisis de seguridad
	\chapter{Análisis del problema}
 
Queremos reducir al máximo la inseguridad de todas las personas en la vía pública.
Las propuestas que hemos estudiado hasta el momento se concentran en la pedida de auxilio.

Podemos identificar tres niveles dentro del problema:
\begin{enumerate}
  \item \textbf{Sospecha, inseguridad:} No lo sabes con certeza, pero algo te pone alerta, podrías estar en peligro y te gustaría estar acompañado.
  \item \textbf{Miedo:} Estás en una situación totalmente incómoda, no sabes si va a llegar a más, pero necesitas compañía.
  \item \textbf{Pánico:} Es necesaria una ayuda inmediata y sería pertinente avisar a los cuerpos de seguridad.
\end{enumerate}

Para los dos primeros, que a su vez son las mas comunes, no hay propuestas.

	% Desarrollo bajo sprints: 
	% 	1. Permitir registros y login de usuarios
	% 	2. Desarrollo del sistema de incidencias
	% 	3. Desarrollo del sistema de denuncias administrativas y accidentes
	% 	4. Desarrollo del sistema de croquis
	%   5. Instalación de la aplicación de manera automática
	\chapter{Implementación}

La implementación del software se ha dividido en hitos. Estos, han sido definidos en Github
y cada uno de ellos contiene un grupo de \textit{issues} que se corresponden con las distintas
mejoras que se han ido incorporando al software a lo largo de su desarrollo.\\



	% Presupuesto

	% Conclusiones
	\chapter{Conclusiones y trabajos futuros}



	% Trabajos futuros


	
	\newpage
	\bibliography{bibliografia}
	\bibliographystyle{plain}
	
\end{document}

